\chapter{Wykład VIII}

\section{Ciąg kompozycyjny. Nilpotentna. Prosta.}



\begin{definition}[Ciąg subnormalny]
	Skończony ciąg podgrup $G= G= \unrhd G_1\unrhd \ldots\unrhd G_{n}= \{e\} $ 

	nazywamy ciągiem subnormalnym $G$. I każdy $\mfaktor{G_{i}}{G_{i+1}} $ faktorami (tego ciągu)

	Gdy wszystki $G_{i}$ są normalne w G, to ciąg nazywamy normalnym.
\end{definition} 
\begin{remark}
	\begin{enumerate}[label=\arabic*)]
		\item Może się zdarzyć, że $G\unrhd G_1\unrhd $, ale $G \unrhd G_2$, np, $A_4\unrhd A'_{4} \unrhd \{ \id, (1,2)(3,4) \} $
		\item $G$ jest rozwiązalna $\iff $ ciąg komutantów jest $\left[ sub \right] $-normalny
		\item $\mfaktor{G^{(i)}}{G^{(i+1)}} = \left( G^{(i)} \right) ^{ab}$
	\end{enumerate}
\end{remark} 
\begin{theorem}
	Następujące warunki są równoważne
	\begin{enumerate}[label=\arabic*)]
		\item $G$ jest rozwiązalna
		\item Istnieje ciąg normalny w $G$- faktorów abelowych
		\item Istnieje ciąg sub normalny w $G$ o faktorach abelowych
	\end{enumerate}
\end{theorem} 
\begin{proof}
	TODO
\end{proof} 

\begin{lemma}
	\begin{enumerate}[label=\alph*)]
		\item Podgrupa grup rozwiązalnych jest rozwiązalna
		\item Obraz grupy rozwiązalnej przez homomorfizm jest grupą rozwiązalną
	\end{enumerate}
\end{lemma} 
\begin{proof}
	TODO
\end{proof} 
\begin{theorem}
	Załóżmy, że $N\unlhd G$
	
	Wtedy $G$ jest rozwiązalna $\iff$ $N$ i $\mfaktor{G}{N} $ jest rozwiązalna
\end{theorem} 
\begin{proof}
	TODO
\end{proof} 
\begin{theorem}
	Każda $p$-grupa jest rozwiązalna
\end{theorem} 
\begin{proof}
	TODO
\end{proof} 

\textbf{Definicja i Twierdzenie w jednym}
\begin{enumerate}[label=\arabic*.]
	\item Niech $G_{(0) := G}$, $G_{(n+1)} := \left[ G, G_{(n)} \right] $. Jest to dolny ciąg centralny.
Mówimy że $G$ jest \textbf{nilpotentna} gdy $\exists\left( n\in \N \right) \left( G_{(b)} = \{e\}  \right) $ 

Ponieważ $\forall\left( n \right) \left( G^{(n)} \subseteq G_{(n)} \right) $ czyli 

$G\text{ abelowa}\implies G\text{ nilpotentna}\implies G\text{ rozwiązalna}$ 

Do pkazania:

\begin{enumerate}[label=\roman*)]
	\item Każda $p$-grupa jest nilpotentna
	\item Przykładowo $S_{3}$ jest rozwiązalna ale nie jest nilpotentna (dowód)
\end{enumerate}

 	\item Definiujemy $Z_0 := \{e\}, Z_1:= Z(G), Z_{n+1} :=  \prod_{n=1} \left[ Z( \mfaktor{G}{Z_{n}}  \right], \prod_{n=1}: \overset{\text{iloraz}}{\to} \mfaktor{G}{Z_{n}}   $
	Wtedy $\{e\} \unlhd Z_0\unlhd Z_1\unlhd \ldots$ jest to \textbf{główny ciąg centralny} 
\end{enumerate}

Do pokazania:
$G$ jest nilpotentna $\iff \exists\left( n \right)\left( Z_{n}= G \right)  $ 

\textbf{Przykłądy}

\begin{enumerate}[label=\arabic*)]
	\item $\left| G \right| p\cdot q$, p i q są pierwsze to $G$ jest rozwiązalna
	\item $\left| G \right| = 110 \implies G$ rozwiązalna.
\end{enumerate}
TUTaj dowody tych przykładów.

\subsection{Grupa prosta}

\begin{definition}
	Mówimy że $G$ jest grupą prostą, gdy $G$ jest nietrywialna oraz
	\[
	\forall\left( N\unlhd G \right)\left( N = \{e\} \text{ lub } N=G \right)  
	.\] 
	$N$ jest podgrupą niewłaściwą(czyli?)
	
\end{definition} 
\textbf{Przykłady} 
\begin{enumerate}[label=\arabic*)]
	\item $\mathbb{Z} _{p}$ p jest liczbą pierwzą
	\item $A_{n}$ 
	\item $PSL_{n}(\mathbb{R} ):= \mfaktor{SL_{n}(\mathbb{R} )}{Z(SL_{n}(\mathbb{R} )} $ 
	\item $PSL_{n}(\mathbb{Z} _{p})$
\end{enumerate}

Zmierzamy do odpowiedzi na pytanie na ile dowolną grupę skończoną można opisać przez grupy proste.

\begin{definition}
	Niech
	\[
	G = H_0 \unrhd H_1 \unrhd \ldots\unrhd H_{n} = \{e\} 
	.\] 
	oraz
	\[
	G = N_0 \unrhd N_1 \unrhd \ldots \unrhd N_{k} = \{e\} 
.\] 
będą ciągami subnormalnymi grupy G.

\begin{enumerate}[label=\arabic*.]
	\item Ciąg $\left( H_{i} \right)_{i}$ nazywamy zagęszczeniem ciągu $\left( N_{i} \right)_{i}$ jeśli
		\[
		\forall\left( j \right) \exists\left( i \right)\left( N_{j}= H_{i} \right)  
		.\] 
	\item Ciąg $\left( H_{i} \right)_{i}, \left( W_{j} \right)_{j} $ są równoważne, gdy $n=k$ oraz
		$\exists\left( \sigma \in S_{k} \right)\left( \forall\left( i \in \{1,\ldots,k\}  \right)  \right)  \left( \mfaktor{H_{i}}{H_{i+1}} \cong \mfaktor{N_{\sigma (i)}}{N_{\sigma (i+1)}}  \right) $
		czyli faktory są iizomorficzne po ich permutacji
\end{enumerate}
\end{definition} 

\begin{theorem}[Scheier]
	Każde dwa ciągi subnormalne $G$ mają równoważne (subnormalne) zagęszczenie.
\end{theorem} 
\begin{proof}
	todo
\end{proof} 
\begin{theorem}
	$A\unlhd B\unlhd 6$ oraz $H\unlhd G$. Wtedy 
	\[
	\mfaktor{BH}{AH} \cong \mfaktor{B}{A(H\cap B)} 
	.\] 
	w szczegolnosci
	$\mfaktor{BH}{H} \cong \mfaktor{B}{H\cap B} $
\end{theorem} 

\begin{definition}
       Ciąg kompozycyjny $G$ to ciąg subnormalny $G$ o faktorach, które są grupami prostymi	
\end{definition} 

\textbf{Przykłądy} 
\begin{enumerate}[label=\arabic*.]
	\item $\{0\} \unlhd \{0,2\} \unlhd \mathbb{Z}_{4}$ 
	\item $\{\id\} \unlhd \{\id, (1,2)(3,4)\}  \unlhd A'_{4} \unlhd A_{4} \unlhd S_4$ 
	\item $\left( Q, + \right) $ nie ma ciągu kompozycyjnego
	\item $\forall\left( n\ge  5 \right) \left( \{\id \} \unlhd A_{n} \unlhd S_{n} \right) $ ciąg kompozycyjny $S_{n}$
\end{enumerate}
\begin{lemma}
	Ciągu kompozycyjnego nie można istotnie zagęscić
\end{lemma} 
\begin{proof}
	TODO
\end{proof} 
\begin{theorem}[Jordana-Koldera]
	Dowolne dwa ciągi kompozycyjne grupy G są równoważne.
\end{theorem} 
\begin{remark}
	Rozważając ciąg kompozycyjny $\mathbb{Z}_{n}$ widać, że tw J-K implikuje jednoznaczność rozkładu n na czynniki pierwsze.
\end{remark} 
\begin{theorem}
	Jeśli $G$ jest skończona to $G$ ma ciąg kompozycyjny.
\end{theorem} 
\begin{remark}
	Ostatnie twierdzenia(które) dają:

	$G$ skonczona implikuje skonczone(proste) faktory ciągu kompozycyjnego G. Są one yznaczone z dokladnoscia do permutajcji jednoznacznie. Czyli mamy rodzaj rozkładu skonczone grupy na grupy proste które dla $G$ skonczonej są sklasyfikowane. Na ile ten fakt klasyfikuje grupy skonczone?
	Nie do konca bo stety i niestety przyporządkowanie 
	\[
	G \mapsto \{G_1,G_2,\ldots, G_{i}\} \text{faktory ciągu kompozycyjnego}
	.\] 
	lub 
	 \[
	G \mapsto \left( G_1,\ldots,G_{i} \right) \text{ ciągu faktorów}
	.\] 
	Nie są 1-1.

\end{remark} 

\textbf{Przykłąd} 
\begin{enumerate}[label=\arabic*)]
	\item $S_3 \unrhd A_3 \unrhd \{\id \} $ ciąg kompozycyjny o faktorach: $\mathbb{Z}_{2}, \mathbb{Z}_{3}$ oraz $\mathbb{Z}_{6} \unrhd \{0,2,4\} \unrhd \{0\}$ ciąg kompozycyjny o faktorach $\mathbb{Z} _{6}$ 

		Ale $S_3\cong \mathbb{Z}_{6}$ 

		W tym przypadku mamy akurat $\mathbb{Z} _{6}\cong \mathbb{Z} _{3}\times \mathbb{Z} _{2}$ oraz $S_3 \cong \mathbb{Z} _{3}\times \mathbb{Z} _{2}$ czyli mamy tzw rozkład półprosty na faktory
	\item $\mathbb{Z} _{4} \unrhd \{0,2\} \unrhd \{0\} $ ciąg kompzycyjne $\mathbb{Z}_{4}$ o faktorach $\mathbb{Z}_{2}, \mathbb{Z} _{2}$.

		$\mathbb{Z} _{2}\times \mathbb{Z} _{2}\unrhd \mathbb{Z} _{2} \times \{0\} \unrhd \left[ (0,0) \right]  $
		 ciąg kompozycyjny $\mathbb{Z}_{2}\times \mathbb{Z} _{2}$ o faktorach $\mathbb{Z} _{2},\mathbb{Z} _{4}$ 
		 ake $\mathbb{Z}_{4}$ nie musi byc produktem półprostym $\mathbb{Z}_{2} $ i $\mathbb{Z} _{2}$
\end{enumerate}
